% !TeX root = ../VCL note.tex
\section{VCL基础知识和技能}

\subsection{颜色，颜色感知和可视化}

人感知颜色靠的是对不同波长的光敏感的三种视锥细胞，这提供了颜色离散表达的基础。三种感光开展出RGB三色系统。颜色的匹配函数如下，是一个积分的办法。

\[
R=\int_0^{\infty}I(\lambda)\bar{r}(\lambda)d\lambda
\]
\[
G=\int_0^{\infty}I(\lambda)\bar{g}(\lambda)d\lambda
\]
\[
B=\int_0^{\infty}I(\lambda)\bar{b}(\lambda)d\lambda
\]

三种通道和原色的光由离散合成有两种机制：基于发光叠加相加混光（显示器发光原理）和基于白光吸收的相减混光（涂料和打印原理）。

计算机视觉基于人对光的感知。在人的感知之下，红色有前进的效应，蓝色会有看起来后退的效应。此即所谓色彩立体视觉。同时在人的认知下颜色也有冷色和暖色之分（这真的是课堂PPT里面的）。颜色信息在人眼中具有恒常性，不论在何种光照之下人对于颜色的认知不会发生大的改变。

当然这种逻辑感知也会带来相应的视觉错觉，譬如色诱导，使用补色可以让人自动脑补对应的颜色；再比如说边缘效应和移动错觉（动态立体视觉）。

人眼对于不同的光敏感性不同，对不同波段的光，施加灰度和模糊处理可以得到不同的认知效果（灰度指的是介于黑白之间的明暗程度）。

